\manualchapter{Control reference}{sec:controls}
Panel map: \textbf{1} pitch and pitch/FM inputs;
\textbf{2} FM and level;
\textbf{3} noise mode and CV;
\textbf{4} audio outputs;
\textbf{5} shared clock, filter and LFSR switches/CV.

The four voice columns repeat frequency, FM, noise and level controls above
their sockets and outputs. The right-hand switches and CV inputs affect shared
clocking, filtering and the shift register.

\begin{tabularx}{\textwidth}{@{}l l X@{}}
\toprule Control & Range; default & Operation \\\midrule
Frequency & $\pm2.5$ oct.; 0 & C4 reference in the calculation; V/OCT adds octaves. \\
FM depth & $-1$--$+1$; 0 & Adds $aV/5$ octaves. \\
Noise & 0--7; 7 & Adds one setting per volt, truncates and clamps. \\
Level & 0--15; 7 & Knob times $V/10$, rounded and clamped. \\
\bottomrule
\end{tabularx}

Pitch, FM, Noise and Level normal forward independently, starting at 0, 5,
0 and 10\,V. Noise selects register-bit combinations, not a smooth wet/dry
mix. Its setting can change whether the clock produces a periodic tone or
a shift-register pattern.

The global switches default off. Low Frequency changes the shared clock;
Tone 1 and Tone 3 High Frequency select faster clocks for those channels.
High-Pass Tone 1 from Tone 3 and Tone 2 from Tone 4 couple the named channels.
LFSR changes the shared noise sequence. Their CV inputs switch at 2\,V high
and 0.01\,V low, reversing the panel state while high. The hardware's paired
16-bit modes are not exposed by this module.
